\clearpage
\VerdierPhysicalPageAnchor{0136}%
\begingroup
\small
\setlength{\emergencystretch}{1.5em}%
\setlength{\abovedisplayskip}{1pt plus 1pt minus 1pt}%
\setlength{\belowdisplayskip}{1pt plus 1pt minus 1pt}%

\VerdierSemanticAnchor{ch2-2-2-1-fl-source-target-maps}
\noindent 以 $\mathbf{s}$ 与 $\mathbf{b}\colon
\operatorname{Fl}(\mathcal C(S^{-1}))\rightrightarrows
\operatorname{Ob}(\mathcal C)$ 表示映射：

\VerdierSemanticAnchor{ch2-2-2-1-source-target-formula}
\[
  m\in\operatorname{Hom}_{\mathcal C(S^{-1})}(X,Y),\qquad
  \mathbf{s}(m)=X,\qquad \mathbf{b}(m)=Y,
\]

\VerdierSemanticAnchor{ch2-2-2-1-identity-map}
\noindent 并以 $\mathbf{id}\colon\operatorname{Ob}(\mathcal C)
\to\operatorname{Fl}(\mathcal C(S^{-1}))$ 表示映射
$X\mapsto Q(\mathbf{id}_X)$。
\VerdierSemanticAnchor{ch2-2-2-1-category-diagram}%
由此定义图式
$\mathcal C(S^{-1})=(\operatorname{Ob}(\mathcal C),
\operatorname{Fl}(\mathcal C(S^{-1})),\mathbf{s},\mathbf{b},\mathbf{id})$。
此外还有一对滥用记号表示为下式的映射：

\VerdierSemanticAnchor{ch2-2-2-1-q-pair}
\[
  Q=\bigl(\mathbf{id}_{\operatorname{Ob}(\mathcal C)}\colon
  \operatorname{Ob}(\mathcal C)\to\operatorname{Ob}(\mathcal C),\;
  Q\colon\operatorname{Fl}(\mathcal C)\to
  \operatorname{Fl}(\mathcal C(S^{-1}))\bigr)\,.
\]

\VerdierSemanticAnchor{ch2-2-2-2}
\smallskip
\noindent\textbf{定理 2.2.2.}\quad
\begingroup\itshape
\VerdierSemanticAnchor{ch2-2-2-2-a}%
a) 在 $\operatorname{Fl}(\mathcal C(S^{-1}))$ 上存在唯一的部分定义的
复合法则，使得：

\VerdierSemanticAnchor{ch2-2-2-2-a-1}
\par\noindent\hspace*{2em}1) $\mathcal C(S^{-1})$ 是一个范畴。

\VerdierSemanticAnchor{ch2-2-2-2-a-2}
\par\noindent\hspace*{2em}2) $Q(s)$ 是同构，且 $Q(s)^{-1}$ 是其逆
$(s\in S)$。

\VerdierSemanticAnchor{ch2-2-2-2-a-3}
\par\noindent\hspace*{2em}3) 对每个图式
$D=X\xleftarrow{s}X'\xrightarrow{m}Y$，$s\in S$，都有
$\overline D=Q(m)Q(s)^{-1}$
$(\overline D\in\operatorname{Hom}_{\mathcal C(S^{-1})}(X,Y)$ 是
$D$ 的等价类）。

\VerdierSemanticAnchor{ch2-2-2-2-a-4}
\par\noindent\hspace*{2em}4) $Q$ 是一个函子。

\VerdierSemanticAnchor{ch2-2-2-2-b}
\par\noindent b) $\mathcal C(S^{-1})$ 的每个态射都可写成复合
$Q(t)^{-1}Q(m)$，其中 $m\in\operatorname{Fl}(\mathcal C)$、$t\in S$
（分别地，写成 $Q(p)Q(s)^{-1}$，其中
$p\in\operatorname{Fl}(\mathcal C)$、$s\in S$）。

\VerdierSemanticAnchor{ch2-2-2-2-c}
\par\noindent c) 函子 $Q$ 与有限射影极限及有限归纳极限可交换。

\VerdierSemanticAnchor{ch2-2-2-2-d}
\par\noindent d) 对每个范畴 $\mathcal A$，与函子 $Q$ 的复合定义函子：

\VerdierSemanticAnchor{ch2-2-2-2-hom-q-a}
\[
  \mathcal{H}\!om(Q,\mathcal A)\colon
  \mathcal{H}\!om(\mathcal C(S^{-1}),\mathcal A)
  \longrightarrow \mathcal{H}\!om(\mathcal C,\mathcal A)
\]

\noindent\textup{（$\mathcal{H}\!om$ 表示函子范畴）。}
\VerdierSemanticAnchor{ch2-2-2-2-universal-image}%
函子 $\mathcal{H}\!om(Q,\mathcal A)$ 在对象上单射且充分忠实；
$\mathcal{H}\!om(\mathcal C(S^{-1}),\mathcal A)$ 的像，是
$\mathcal{H}\!om(\mathcal C,\mathcal A)$ 的严格满子范畴，其对象是
那些把 $S$ 中的态射送到同构的从 $\mathcal C$ 到 $\mathcal A$ 的函子。
\par\endgroup

\VerdierSemanticAnchor{ch2-2-2-2-proof-omitted}
该定理的证明留给读者；读者可参阅 [16]。
\VerdierSemanticAnchor{ch2-2-2-2-composition-setup}%
这里只描述 $\mathcal C(S^{-1})$ 中态射的复合。设
$\overline D_1\in\operatorname{Hom}_{\mathcal C(S^{-1})}(X,Y)$ 与
$\overline D_2\in\operatorname{Hom}_{\mathcal C(S^{-1})}(Y,Z)$ 为
$\mathcal C(S^{-1})$ 中两个可复合的态射。

\VerdierSemanticAnchor{ch2-2-2-2-roof-pair}
设 $D_1=X\xleftarrow{s}X'\xrightarrow{m}Y$ 与
$D_2=Y\xleftarrow{t}Y'\xrightarrow{p}Z$ 为 $\mathcal C$ 中分别属于
$\overline D_1$ 与 $\overline D_2$ 的两个图式。
\VerdierSemanticAnchor{ch2-2-2-2-sm2-completion-input}%
利用 SM2)，把图式

\begin{center}
  \begin{tikzpicture}[>=stealth,line width=.38pt,x=1.42cm,y=.90cm,
    scale=.92,transform shape]
    \node (Xp) at (-1.20,1.10) {$X'$};
    \node (Yp) at ( 1.20,1.10) {$Y'$};
    \node (Y)  at ( 0,-.70) {$Y$};
    \draw[->] (Xp) -- node[left=2pt] {$m$} (Y);
    \draw[->] (Yp) -- node[right=2pt] {$t$} (Y);
    \node at (3.00,.25) {$t\in S$};
  \end{tikzpicture}
\end{center}

\VerdierSemanticAnchor{ch2-2-2-2-p137-boundary}
\endgroup
